\chapter{检索服务表示}

\section{规范表示与派生视图}

权威数据保存认知对象及修订的语义身份。检索服务可为同一修订建立规范表示、语义片段、有来源约束的召回线索、实体/标签倒排项、词项文档和稠密向量视图。视图粒度和数量不会改变认知身份。\cite{nous-research-package}

每个稠密向量视图绑定完整的向量空间标识，包括模型/权重、任务、分词器、预处理、池化、归一化和维度。不同空间分别检索和排序；跨空间合并须通过明确的转换或分区规则。

\section{可独立重建的投影类别}

精确索引、词项索引、稠密向量、拓扑和摘要视图可分别失效、重建和发布。权威数据变更与投影水位位于事务路径；每个代次记录输入水位、实现标识、配置摘要、制品位置/校验和及向量空间。构建区通过验证后，再原子切换当前指针。\cite{nous-research-package}

当前 AuthorityStore 已按投影类别记录水位并发布独立代次；实际服务使用仓库配置的本地 Tantivy、USearch、精确倒排和 petgraph 路径。这些是当前工程选型，后续按目标场景评测调整。\cite{nous-code-snapshot}

\section{稠密检索后端仍待比较}

研究资料包把 Qdrant 和 pgvector 列为本机稠密检索候选。Qdrant 适合命名向量、多向量及独立服务生命周期；pgvector 可减少独立服务进程，并与 PostgreSQL 过滤和备份运维结合。项目组将依据目标硬件、多视图查询、过滤选择性、重建、崩溃恢复和备份恢复进行对照评测，再确定默认方案。\cite{nous-research-package}

\section{投影不拥有认知真值}

认知身份由权威对象和修订标识；向量点与索引项属于服务视图。权威数据维护长期认知，服务代次可独立删除和重建；中心性、传播能量、向量得分和重排结果仅用于检索排序。
